\documentclass[10pt,a4paper]{article}
\usepackage[margin=20mm,headheight=15pt,headsep=7mm,footskip=10mm]{geometry}
\usepackage{fontspec,amsmath,amssymb,mathtools,microtype}
\setmainfont{lmroman10-regular.otf}[BoldFont=lmroman10-bold.otf,ItalicFont=lmroman10-italic.otf,BoldItalicFont=lmroman10-bolditalic.otf]
\setsansfont{lmsans10-regular.otf}[BoldFont=lmsans10-bold.otf,ItalicFont=lmsans10-oblique.otf,BoldItalicFont=lmsans10-boldoblique.otf]
\setmonofont[Scale=.86]{lmmono10-regular.otf}[ItalicFont=lmmono10-italic.otf]
\usepackage{ragged2e}\usepackage{xcolor,booktabs,longtable,array,graphicx,enumitem,fancyhdr,needspace,titlesec,hyperref,xurl}
\definecolor{AtlasNavy}{HTML}{18334C}\definecolor{AtlasTeal}{HTML}{116D73}\definecolor{AtlasGray}{HTML}{52616D}
\hypersetup{colorlinks=true,linkcolor=AtlasTeal,urlcolor=AtlasTeal,pdfcreator={Complexity Atlas},unicode=true}
\urlstyle{same}\setlength{\parindent}{0pt}\setlength{\parskip}{4.8pt}\setlength{\emergencystretch}{3em}
\renewcommand{\baselinestretch}{1.08}\setcounter{secnumdepth}{0}\setlist{nosep,leftmargin=1.5em}
\titleformat{\section}{\Large\bfseries\sffamily\color{AtlasNavy}\raggedright}{}{0em}{}
\titleformat{\subsection}{\large\bfseries\sffamily\color{AtlasTeal}\raggedright}{}{0em}{}
\titleformat{\subsubsection}{\normalsize\bfseries\color{AtlasNavy}\raggedright}{}{0em}{}
\titlespacing*{\section}{0pt}{14pt}{6pt}\titlespacing*{\subsection}{0pt}{10pt}{5pt}
\pagestyle{fancy}\fancyhf{}\fancyhead[L]{\small\sffamily\color{AtlasGray}COMPLEXITY ATLAS}
\fancyhead[R]{\small\sffamily\color{AtlasGray}ENGLISH EDITION}\fancyfoot[C]{\small\thepage}\renewcommand{\headrulewidth}{.3pt}
\newcommand{\AtlasMath}[1]{\par\begingroup\centering\setbox0=\hbox{$\displaystyle #1$}\ifdim\wd0>1.4\linewidth\typeout{ATLAS WIDE FORMULA: \the\wd0}\fi\ifdim\wd0>\linewidth\resizebox{\linewidth}{!}{\box0}\else\box0\fi\par\endgroup}
\newcommand{\AtlasInlineMath}[1]{\begingroup\setbox0=\hbox{$#1$}\ifdim\wd0>.4\linewidth\endgroup\AtlasMath{#1}\else\endgroup$#1$\fi}
\pdfstringdefDisableCommands{\def\AtlasInlineMath#1{#1}}
\DeclareMathOperator{\polylog}{polylog}\DeclareMathOperator{\dist}{dist}\DeclareMathOperator{\diam}{diam}
\DeclareMathOperator{\rank}{rank}\DeclareMathOperator{\tr}{tr}\DeclareMathOperator{\diag}{diag}
\DeclareMathOperator{\prox}{prox}\DeclareMathOperator{\KL}{KL}\DeclareMathOperator{\Var}{Var}
\providecommand{\eps}{\varepsilon}\providecommand{\R}{\mathbb R}\providecommand{\E}{\mathbb E}
\providecommand{\mathbbm}[1]{\mathbf{#1}}\providecommand{\mathds}[1]{\mathbf{#1}}
\providecommand{\llbracket}{[\![}\providecommand{\rrbracket}{]\!]}
\begin{document}
\raggedright

{\LARGE\bfseries\sffamily\color{AtlasNavy} Nonconvex-strongly concave optimization with reusable samples and mean-square smoothness.\par}\vspace{5pt}
\textbf{Setting ID.} \texttt{CV\_\allowbreak{}S\_\allowbreak{}NCSC\_\allowbreak{}MSS\_\allowbreak{}REUSE}

Minimax optimization

Solve \AtlasInlineMath{\min_x\Phi(x)}, where \AtlasInlineMath{\Phi(x)=\max_y f(x,y)} and \AtlasInlineMath{f=\mathbb E_\xi f_\xi}.

\textbf{Status.} Bounds missing

\textbf{TASK tags.} task family: Minimax

\textbf{MODEL tags.} convexity: Nonconvex strongly concave; convexity: Strong concavity; initialization: Initialization; smoothness: Sample mean square

\textbf{ORACLE tags.} feedback quality: Exact; feedback quality: Stochastic unbiased bounded variance; interface family: Reusable sample gradient; operation: Sample draw; operation: Sample gradient; public information: Hidden objects explicitly declared; public information: Public inputs explicitly declared; scalar derivative order: Scalar-function first derivative

\textbf{ALGORITHMS tags.} adaptivity: Adaptive; randomization: Private randomness allowed; restriction: Explicitly no additional algorithm restrictions; restriction: No additional span restriction explicitly imposed; restriction: No additional zero-respecting restriction explicitly imposed

\textbf{OUTPUT tags.} certificate: No additional output certificate required; object: Point

\textbf{METRIC tags.} criterion: Reduced-objective gradient norm

\textbf{GUARANTEE tags.} aggregation: Per instance; aggregation: Second moment of the complete error; mode: Expectation; quantifier: Algorithm before instance; randomness: Algorithm private seed; randomness: Oracle samples

\textbf{RESOURCE tags.} accounting: Joint coordinates retained; budget: Fixed worst-case budget; measure: N independent sample; measure: N sample gradient; setup: Initialization included

\needspace{7\baselineskip}\section{1. Setting specification}
\needspace{4\baselineskip}\subsection{Task}
Solve \AtlasInlineMath{\min_x\Phi(x)}, where \AtlasInlineMath{\Phi(x)=\max_y f(x,y)} and \AtlasInlineMath{f=\mathbb E_\xi f_\xi}.

\needspace{4\baselineskip}\subsection{Admissible model}
\AtlasInlineMath{f:\mathbb{R}^{d_x}\times\mathbb{R}^{d_y}\to\mathbb{R}}, with both blocks unconstrained and Euclidean norms throughout. The joint gradient satisfies \AtlasInlineMath{\|\nabla f(z)-\nabla f(w)\|_2\le \beta\|z-w\|_2}.

For every \AtlasInlineMath{x}, \AtlasInlineMath{f(x,\cdot)} is uniformly \AtlasInlineMath{\mu_y}-strongly concave. Thus \AtlasInlineMath{y^\star(x)=\arg\max_y f(x,y)} is unique and \AtlasInlineMath{\Phi(x)=f(x,y^\star(x))} is differentiable. Neither \AtlasInlineMath{f(\cdot,y)} nor \AtlasInlineMath{\Phi} is assumed convex.

Assume \AtlasInlineMath{\inf_x\Phi(x)>-\infty}, with an initial point satisfying \AtlasInlineMath{\Phi(x_0)-\inf_x\Phi(x)\le\Delta}.

The initial point additionally satisfies \AtlasInlineMath{\|\nabla_yf(x_0,y_0)\|_2\le G_0}. Every sample function is differentiable. For any fixed points \AtlasInlineMath{z,w}:


\AtlasMath{\mathbb{E}\nabla f_\xi(z)=\nabla f(z),\qquad\mathbb{E}\|\nabla f_\xi(z)-\nabla f(z)\|_2^2\le\sigma^2,}



\AtlasMath{\mathbb{E}\|\nabla f_\xi(z)-\nabla f_\xi(w)\|_2^2\le\beta_{\rm ms}^2\|z-w\|_2^2.}


\needspace{4\baselineskip}\subsection{Oracle interface}
Public inputs include the block dimensions \AtlasInlineMath{d_x,d_y} and \AtlasInlineMath{\Delta}, \AtlasInlineMath{G_0}, \AtlasInlineMath{\beta}, \AtlasInlineMath{\beta_{\rm ms}}, \AtlasInlineMath{\varepsilon}, \AtlasInlineMath{\mu_y}, and \AtlasInlineMath{\sigma}. The stated initial points and feasible-set information are also supplied. An analytic representation of the hidden function, operator, or matrix is not revealed.

Sample identifiers \AtlasInlineMath{\xi} are drawn independently from an unknown distribution fixed before interaction. Identifiers may be stored and used to evaluate \AtlasInlineMath{\nabla f_\xi} at several points in the full space. Every sample-point evaluation is charged.

A paired-sample difference at two points requires two evaluations. Values and Hessians are unavailable. Adaptive points chosen after obtaining a reusable sample identifier need not be unbiased conditional on the entire history; fixed-point unbiasedness and persistent sample coupling retain their model definitions.

\needspace{4\baselineskip}\subsection{Output}
Return \AtlasInlineMath{\widehat x\in\mathbb R^{d_x}}.

The output may be any measurable postprocessing of the history, including an unqueried point, average, or convex combination. It must satisfy the stated domain and feasibility requirements; no additional function value or computable certificate is required.

\needspace{4\baselineskip}\subsection{Loss or accuracy criterion}

\AtlasMath{\|\nabla\Phi(\widehat{x})\|_2}


\needspace{4\baselineskip}\subsection{Resource accounting}
The primary cost is \AtlasInlineMath{N_{\rm sample\text{-}gradient}}, with independent sample count recorded separately. Reusing one sample identity at several points does not turn those evaluations into one gradient call.

Use a fixed worst-case budget: calls never exceed the prespecified \AtlasInlineMath{N} on any instance or random execution path. Earlier stopping within the budget is permitted.

Count all information calls for one instance and one requested accuracy, including initialization. Exact-real arithmetic, storage, and solving known models from public inputs and returned data are uncharged. This query-complexity quantity does not certify arithmetic or bit complexity.

\needspace{4\baselineskip}\subsection{Algorithm views and success guarantees}
Arbitrary measurable sequential adaptive randomized algorithms are allowed. The private seed is independent of the hidden instance. Queries, stopping, and output may use the full history, without additional span, zero-respecting, communication, or memory restrictions.

Choose one algorithm, then require the guarantee for every admissible instance and stochastic-feedback mechanism.


\AtlasMath{\mathbb{E}\!\left[\left(\|\nabla\Phi(\widehat{x})\|_2\right)^2\right]\le\varepsilon^2.}


Expectation is over private algorithmic randomness and the random samples actually used.

\textbf{Budget convention.} fixed\_worst\_case\_cap

\needspace{7\baselineskip}\section{2. Notation and parameters}
\AtlasInlineMath{\Delta}: A public bound on the initial reduced-objective value above its infimum. Attainment of the infimum is not required. \textbf{Domain:} Positive real number.

\AtlasInlineMath{\Phi(x_0)-\inf_x\Phi(x)\le\Delta}

\AtlasInlineMath{G_0}: Euclidean norm bound on the dual gradient at the public initial point \textbf{Domain:} A nonnegative real.

\AtlasInlineMath{\|\nabla_yf(x_0,y_0)\|_2\le G_0}

\AtlasInlineMath{\beta}: Euclidean Lipschitz upper bound for the population payoff's complete joint gradient. \textbf{Domain:} Positive real number.

\AtlasInlineMath{\|\nabla f(z)-\nabla f(w)\|_2\le\beta\|z-w\|_2}

\AtlasInlineMath{\beta_{\rm ms}}: Mean-square Lipschitz upper bound for a same-sample gradient difference between two points. \textbf{Domain:} Positive real number.

\AtlasInlineMath{d_x}: Primal-space dimension \textbf{Domain:} Positive integer.

\AtlasInlineMath{d_y}: The dimension of the dual-variable space. \textbf{Domain:} Positive integer.

\AtlasInlineMath{\varepsilon}: Euclidean norm tolerance for the true reduced gradient; the mean-square contract targets its square \AtlasInlineMath{\varepsilon^2}. \textbf{Domain:} Positive real number.

\AtlasInlineMath{\kappa}: Ratio of the population smoothness upper bound, or mean-square average smoothness for finite sums, to the dual strong-concavity lower bound \textbf{Domain:} \AtlasInlineMath{\kappa\ge1}.

\AtlasInlineMath{\kappa=\beta/\mu_y}

\AtlasInlineMath{\mu_y}: A uniform strong-convexity lower bound for -f(x,·) at every fixed x; equivalently, the dual strong-concavity modulus of f. \textbf{Domain:} \AtlasInlineMath{0<\mu_y\le\beta}.

\AtlasInlineMath{f(x,v)\le f(x,y)+\langle\nabla_yf(x,y),v-y\rangle-\frac{\mu_y}{2}\|v-y\|_2^2}

\AtlasInlineMath{\sigma}: Standard-deviation upper bound for noise in the full stochastic response, controlling the joint squared norm rather than summing independent block budgets. \textbf{Domain:} A nonnegative real.

\AtlasInlineMath{\mathbb E_\xi\|\nabla f_\xi(z)-\nabla f(z)\|_2^2\le\sigma^2\quad\text{for every fixed }z}.

\AtlasInlineMath{\|\cdot\|_2,\langle\cdot,\cdot\rangle}: Euclidean norm and inner product; the product-space norm is the square root of the sum of the squared block norms. \textbf{Domain:} The specified finite-dimensional real vector space.

\AtlasInlineMath{\|v\|_2=\sqrt{\langle v,v\rangle}}

\AtlasInlineMath{f}: Unknown real-valued payoff, with primal and dual input blocks. \textbf{Domain:} \AtlasInlineMath{f:(x,y)\mapsto f(x,y)\in\mathbb R}.

\AtlasInlineMath{x,y,z}: Primal variable, dual variable, and joint point \textbf{Domain:} \AtlasInlineMath{x\in\mathbb R^{d_x},\ y\in\mathbb R^{d_y},\ z=(x,y)\in\mathbb R^{d_x+d_y}}.

\AtlasInlineMath{z=(x,y)}

\AtlasInlineMath{\nabla f,\nabla_x f,\nabla_y f}: Complete joint gradient and its two blocks \textbf{Domain:} \AtlasInlineMath{\nabla_x f\in\mathbb R^{d_x},\ \nabla_y f\in\mathbb R^{d_y}}.

\AtlasInlineMath{\nabla f=(\nabla_x f,\nabla_y f)}

\AtlasInlineMath{\Phi}: Reduced objective after global dual maximization; neither its value nor gradient is an additional free interface \textbf{Domain:} Real-valued function on the primal-variable space.

\AtlasInlineMath{\Phi(x)=\max_{y\in \mathbb R^{d_y}}f(x,y)}

\AtlasInlineMath{y^\star(x)}: The unique global dual maximizer at fixed primal x. \textbf{Domain:} \AtlasInlineMath{y^\star(x)\in\mathbb R^{d_y}}.

\AtlasInlineMath{y^\star(x)=\arg\max_{y\in\mathbb R^{d_y}}f(x,y)}

\AtlasInlineMath{x_0,y_0}: The public initial primal and dual variables. \textbf{Domain:} The respective feasible domains. Both equal zero when the origins are explicitly fixed.

\AtlasInlineMath{\widehat x}: The point or variable blocks ultimately returned by the algorithm, corresponding to the hatted variables in the error formula. \textbf{Domain:} A feasible primal point in \AtlasInlineMath{\mathbb R^{d_x}}.

\AtlasInlineMath{\mathbb E,\Pr}: Expectation and probability over the explicitly allowed algorithm and sample randomness, with the hidden instance fixed. \textbf{Domain:} Operations on the corresponding probability space.

\AtlasInlineMath{\xi,f_\xi}: A retained sample identifier and its fixed differentiable sample function, drawn independently at creation. \textbf{Domain:} \AtlasInlineMath{f_\xi:\mathbb R^{d_x+d_y}\to\mathbb R}.

\AtlasInlineMath{f(z)=\mathbb E_\xi f_\xi(z)}

\AtlasInlineMath{N_{\rm independent}}: Number of independent sample identifiers drawn \textbf{Domain:} Nonnegative integer count; expected work is its expectation.

\AtlasInlineMath{N_{\rm sample\text{-}gradient}}: Gradient evaluations of one sample identifier at one point. \textbf{Domain:} Nonnegative integer count; expected work is its expectation.

\AtlasInlineMath{\mathcal A,I}: An algorithm selected from public inputs, a hidden admissible object fixed before interaction, and an oracle law satisfying the declared conditional means, variances, and independence protocol. \textbf{Domain:} The specified algorithm and model classes.

\AtlasInlineMath{\omega}: A complete realization of algorithmic and sample randomness. \textbf{Domain:} Corresponding sample space.

\AtlasInlineMath{B_{\rm independent},B_{\rm sample\text{-}gradient}}: One algorithm simultaneously satisfies the separate caps on independent sample draws and sample-gradient evaluations. \textbf{Domain:} A pair of nonnegative integers. Sample identifiers may be stored persistently and reused.

\needspace{12\baselineskip}\section{3. Complexity bracket}
\needspace{4\baselineskip}\subsection{Compared quantities}
\needspace{4\baselineskip}\subsubsection{Q1 - target accuracy cost}
\AtlasMath{C^\star_{v,\pi}(\varepsilon;\theta)=\inf\{B\ge0:\exists\mathcal A\in\mathfrak A_v,\ C_\pi(\mathcal A,I)\le B\ \text{and}\ G_v(\mathcal A,I;\varepsilon)\ \forall I\}}
\textbf{Definition.} Minimum charged resource attaining the stated success guarantee

\textbf{Scope.} The full documented parameter domain; no upper or lower bound is admitted for this comparison.

\textbf{Optimization direction.} minimize

Here v is the selected algorithm view, G\_v its complete success guarantee, I an admissible fixed instance, theta the remaining parameters, and C\_pi the declared resource projection. C\_pi uses the view's fixed-path or expected-work budget convention, including initialization. No change of guarantee or scalarization is implicit.

\needspace{4\baselineskip}\subsection{R1 - The declared parameter domain.}
\textbf{Region.} All admissible parameters in the eight-part specification; subregions with two-sided evidence have not yet been separated.

\textbf{Quantity.} Q1

\textbf{Resource.} The primary charged resource specified in the common resource contract.

\textbf{Status.} Both bounds missing

\textbf{Lower bound.} No admitted bound for this comparison

\textbf{Upper bound.} No admitted bound for this comparison

\needspace{7\baselineskip}\section{4. Bounds and supporting evidence}
\end{document}
